\chapter{INTRODUCTION} 

\section{Background of the Study}

Mental health support has become an increasingly important global concern as psychological distress, anxiety, depression, and related conditions continue to affect people's daily functioning, relationships, work, and quality of life. The World Health Organization reported that more than one billion people are living with mental health disorders, with anxiety and depression among the most common conditions worldwide \cite{who2025mentalhealth}. Despite increasing awareness of mental health issues, access to professional psychological care remains limited, particularly in low- and middle-income countries where treatment gaps are still high \cite{who2025mentalhealth,kazdin2019mentalhealth}. These challenges have encouraged researchers to explore scalable forms of psychological support, including digital mental health interventions and large language model-based counseling systems. Recent studies suggest that large language models may support mental health-related tasks, but their use in counseling requires careful evaluation because therapeutic dialogue involves empathy, safety, contextual understanding, and professional boundaries \cite{hua2025scopingreview,nguyen2025llmcounseling}.

Large language models have shown potential for producing fluent and empathetic responses that fit counseling-like contexts. However, therapeutic communication is not only a matter of generating a warm or natural reply. A helpful response also needs to preserve trust, respect the client's autonomy, and support a collaborative relationship between the client and the therapist. These qualities are central to client-centered counseling and therapeutic alliance, where empathy, agreement on goals, and agreement on therapeutic tasks influence the quality of the interaction \cite{rogers1957necessary,bordin1979alliance}. For LLM-based counseling systems, this means that response generation should be grounded not only in language quality, but also in therapeutic principles and interactional sensitivity.

This thesis focuses on simulated therapeutic dialogue involving resistance and ambivalence. Motivational Interviewing (MI) and Cognitive Behavioral Therapy (CBT) provide two complementary foundations for this setting. MI emphasizes collaboration, autonomy, reflective listening, affirmation, and eliciting the client's own reasons for change, making it useful when the client is uncertain or not yet ready to engage deeply with change \cite{miller2013motivational,bischof2021motivational}. CBT provides structured methods for identifying unhelpful thoughts, examining evidence, developing more balanced interpretations, and supporting behavioral change \cite{beck2020cognitive,chand2023cognitive}. However, CBT's structured and problem-focused nature can become a limitation when the client is resistant or not yet ready for direct cognitive analysis. If introduced too early, CBT-style questioning may feel directive, reduce the client's sense of autonomy, weaken the therapeutic alliance, and increase defensiveness. Together, MI and CBT offer a useful combination for supporting clients who may move between hesitation, exploration, and readiness for more structured therapeutic work.

The integration of MI and CBT has also been explored in conventional mental health treatment outside the LLM context. Arkowitz and Westra discussed the use of MI with CBT for depression and anxiety, especially where resistance and ambivalence affect treatment engagement \cite{arkowitz2004integrating}. In generalized anxiety disorder, studies have examined MI as a pretreatment or integrated component of CBT, with findings suggesting potential benefits for clients with higher severity, resistance, or ambivalence \cite{westra2009adding,westra2016integrating}. Reviews of MI as an adjunct to CBT for anxiety disorders also suggest that MI-CBT can support treatment initiation, engagement, and clinical outcomes, although the literature still points to the importance of how and when MI is integrated with CBT \cite{randall2017motivational,marker2018efficacy}. This clinical background supports the motivation of this thesis: MI can improve engagement with CBT, but MI-CBT integration also creates an important coordination problem. The challenge is not only combining MI and CBT, but deciding when to remain in MI, when to shift toward CBT, and when to transition gradually between the two during therapeutic dialogue.

The difficulty of coordinating MI and CBT can be understood through the client's readiness for change. The Transtheoretical Model describes behavior change as a staged process that includes precontemplation, contemplation, preparation, action, and maintenance \cite{prochaska1997transtheoretical,raihan2023stages}. A client in precontemplation may minimize the problem or resist direct analysis, while a client in preparation may be more open to structured reflection, cognitive restructuring, or planning. Therefore, the same therapeutic move may be helpful in one stage but ineffective or even counterproductive in another. For example, introducing CBT too early may feel directive and increase resistance, whereas remaining only in MI after the client becomes ready may delay more structured therapeutic progress \cite{arkowitz2004integrating,westra2009adding,westra2016integrating,randall2017motivational,marker2018efficacy}. In this thesis, resistance and readiness are treated as dynamic interaction states that shape whether an MI-style response, a CBT-style response, or a gradual transition between the two is more suitable.

These considerations motivate the need for LLM-based counseling systems that are sensitive not only to the current user message, but also to the likely direction of the dialogue. In therapeutic conversation, the goal is not simply to produce a locally appropriate response, but to choose a response that helps the client move toward greater engagement, openness, collaboration, and readiness. This is especially important in MI-CBT dialogue because different therapeutic moves may produce different near-future client reactions depending on the client's resistance, ambivalence, and stage of change. Therefore, the system needs a mechanism for comparing possible MI-style and CBT-style responses before committing to one response. Short-horizon lookahead provides one way to address this problem, because it allows multiple candidate responses and possible future client reactions to be considered before selecting the final response. Similar search-based ideas have been explored in large language model reasoning, where considering multiple candidate paths can improve decision-making in complex tasks \cite{yao2023tree}. This thesis adapts this general idea to simulated MI-CBT therapeutic dialogue by investigating whether reward-guided lookahead search can support better near-future therapeutic states.

To address this direction, this thesis proposes AIM-RGL, an agentic MI-CBT framework with reward-guided lookahead search for managing client resistance in simulated therapeutic conversations. The framework uses specialist therapist agents to generate candidate responses, including agents for open-ended questioning, reflection, affirmation, summarization, and CBT. A client simulator then produces possible future client reactions based on cognitive conceptualization and stage-of-change information, while a reward model evaluates candidate trajectories using therapeutic dimensions such as alliance, empathy, collaboration, CBT skill, MI quality, and client readiness. Instead of selecting only a locally appropriate response, AIM-RGL aims to select the response expected to support a better near-future therapeutic state. This thesis does not claim to diagnose, treat, replace professional therapists, or evaluate real clinical outcomes; rather, it investigates future-aware response selection in simulated AI-assisted counseling dialogue.

\section{Statement of the Problem}

Although LLM-based counseling systems have become increasingly capable of generating fluent and empathetic responses, many still treat therapeutic response generation mainly as an immediate-context decision. This creates a problem because therapeutic response quality depends not only on how appropriate a response appears in the present turn, but also on how it affects the client's next state. A response may sound supportive but still increase defensiveness, reduce engagement, or weaken the therapeutic alliance if it does not match the client's readiness for change \cite{rogers1957necessary,bordin1979alliance}. Recent evaluations of LLM-based counseling systems also highlight the need to assess these systems using counseling-specific competencies, safety, and therapeutic quality rather than relying only on general language quality \cite{hua2025scopingreview,nguyen2025llmcounseling}. Therefore, the central problem addressed in this thesis is the difficulty of selecting therapeutically appropriate responses for resistant or ambivalent clients when response quality depends not only on current-turn suitability, but also on the client's likely next state.

This problem is especially important when an LLM-based system needs to coordinate MI and CBT responses for resistant or ambivalent clients. MI and CBT support different therapeutic needs: MI helps the client explore ambivalence and motivation, while CBT supports structured examination of thoughts and behavior change \cite{miller2013motivational,bischof2021motivational,beck2020cognitive,chand2023cognitive}. However, the usefulness of each approach depends on when it is applied. If CBT is introduced before the client is ready, it may feel directive and increase resistance; if the system remains only in MI after the client becomes ready, it may miss opportunities for structured therapeutic progress \cite{prochaska1997transtheoretical,raihan2023stages}. Therefore, the problem is not simply choosing between MI and CBT, but selecting the therapeutic move that best fits both the client's current state and likely near-future response.

Prior work on LLM-based counseling has studied memory, stage awareness, strategy selection, and therapeutic response evaluation. For example, ChatThero uses a multi-session, stressor-aware language agent for recovery support, while Persona Drift explores adaptive flow in stage-aware CBT chatbots \cite{wang2025chatthero,yun2025personadrift}. More closely related to this thesis, FLASH proposes a hybrid MI-CBT counseling framework that uses Flashbulb Memory to store high-impact client reactions to specific strategies and dynamically guide later intervention selection based on historical success and failure \cite{anonymous2026flash}. However, these systems mainly adapt through memory, stage awareness, or previously observed interaction signals, and may still have limited ability to compare the immediate downstream effects of multiple candidate therapeutic moves before selecting a final response. This creates a limitation in current-turn response selection, where a response may appear appropriate immediately but may not lead to the most beneficial next client state. This thesis addresses the gap by testing whether reward-guided lookahead search can compare multiple candidate MI-CBT responses and select the response expected to produce a better simulated client trajectory.

\section{Objectives}

The main objective of this thesis is to develop and evaluate AIM-RGL, an agentic MI-CBT framework that uses reward-guided lookahead search to improve therapeutic response selection for managing client resistance in simulated counseling dialogue.

\begin{itemize}[leftmargin=0.5in]
    \item To build a multi-agent therapist architecture that generates candidate responses with specialist agents for open-ended questioning, reflection, affirmation, summarization, and CBT.

    \item To design a simulated client interaction setting grounded in cognitive conceptualization and stage-of-change information, so that client resistance, ambivalence, and readiness can be modeled during therapeutic dialogue.

    \item To train a reward model that estimates the therapeutic value of candidate responses based on dimensions such as therapeutic alliance, empathy, collaboration, CBT skill, MI quality, and client readiness.

    \item To implement reward-guided lookahead search that compares multiple candidate MI-CBT responses and selects the response expected to support a better near-future client trajectory.

    \item To examine how different lookahead horizons, including one-step, two-step, and three-step lookahead, affect the quality of MI-CBT therapeutic response selection.

    \item To evaluate whether AIM-RGL improves therapeutic skill quality, MI-CBT fidelity, therapeutic alliance, change talk, and sustain talk compared with selected baseline counseling systems.

    \item To analyze the effect of different reward model backbones and policy model backbones on the performance of reward-guided MI-CBT response generation.
\end{itemize}

\section{Organization of the Study}

This thesis is organized into five chapters. Chapter 1 introduces the research background, problem statement, objectives, and overall organization of the study. It explains the motivation for developing an LLM-based MI-CBT counseling framework that can manage client resistance through future-aware response selection. Chapter 2 reviews the literature related to therapeutic alliance, client resistance, readiness for change, Motivational Interviewing, Cognitive Behavioral Therapy, LLM-based counseling systems, client simulation, therapeutic evaluation, reward models, and lookahead search. This chapter establishes the clinical and technical foundations for AIM-RGL.

Chapter 3 presents the proposed methodology. It describes the overall AIM-RGL framework, including the multi-agent therapist response generation module, client simulation setting, reward-guided lookahead search mechanism, reward model training, and policy optimization process. Chapter 4 presents the experimental design used to evaluate AIM-RGL. It includes the research questions, hypotheses, independent variables, dependent variables, datasets, model backbones, baseline systems, lookahead configurations, evaluation metrics, and comparison procedures. Chapter 5 concludes the thesis by summarizing the main contributions, discussing the expected implications and limitations of the study, and outlining possible directions for future work.